\documentclass{article}
\usepackage{amsmath, amssymb, amsthm}
\usepackage{hyperref}
\usepackage{geometry}
\geometry{margin=1in}

\title{Erd\H{o}s Problem \#942}
\date{Last updated: 2025-08-31}
\author{Source: erdosproblems.com}

\begin{document}
\maketitle

\noindent\textbf{Status:} OPEN \quad \textbf{No prize}

\noindent\textbf{Tags:} number theory

\medskip

\noindent\textbf{Problem Statement:}

Let $h(n)$ count the number of powerful (if $p\mid m$ then $p^2\mid m$) integers in $[n^2,(n+1)^2)$. Estimate $h(n)$. In particular is there some constant $c&#62;0$ such that\[h(n) &lt; (\log n)^{c+o(1)}\]and, for infinitely many $n$,\[h(n) &gt;(\log n)^{c-o(1)}?\]




Erd\H{o}s writes it is not hard to prove that $\limsup h(n)=\infty$, and that the density $\delta_l$ of integers for which $h(n)=l$ exists and $\sum \delta_l=1$. A proof that $h(n)$ is unbounded is provided by van Doorn in the comments. De Koninck and Luca \cite{DeLu04} have proved, for infinitely many $n$,\[h(n) \gg \left(\frac{\log n}{\log\log n}\right)^{1/3}.\]They also give the density ($\approx 0.275$) of those $n$ such that $h(n)=1$.




References

[DeLu04] De Koninck, Jean-Marie and Luca, Florian, Sur la proximit\'{e} des nombres puissants . Acta Arith. (2004), 149--157.

\noindent\textbf{Related OEIS sequences:} possible


\bigskip
\noindent\small{Source: \url{https://www.erdosproblems.com/942}}
\end{document}
